\section{The Appointed Texts}
\label{sec:appointed-texts}

The Latin below is that of the 1962 typical edition, \work{Missale Romanum}
(Typis Polyglottis Vaticanis), pp.~411--412, nos.~1679--1688, read on the page
images of the CMAA facsimile. Every element's wording already stands in the
Pustet \work{Missale Romanum} of Ratisbon, 1862 (pp.~351--352); the two books
differ only in spelling, pointing, capitalisation, the form of the citations,
the length of one printed conclusion, and the additional orations that the
1862 book prints under the rubrics of its day and the 1962 book no longer
carries. The English of the scriptural elements is the Douay--Rheims Bible in
Bishop Challoner's revision, which translates the Clementine Vulgate at the
Missal's own Vulgate numbering. It is given at the canonical verses, so that
wherever a chant departs from its verse the English and the Latin part company;
each such place is named beneath the text. The English of the Introit's
antiphon and of the three orations is the anonymous translation printed in
\work{The Roman Missal, Translated into the English Language for the Use of the
Laity} (Philadelphia: Eugene Cummiskey, 1861). Neither English is an approved
liturgical translation of the 1962 Missal. Psalms are cited in the Missal's
numbering, with the Hebrew numbering in parentheses.

\subsection{Introit}
\label{proper-introit}

\properlocus{Antiphon without printed citation; verse Ps 77:1 (Heb.\ 78). Missal p.~411, no.~1679.}

\begin{latinproper}
Salus pópuli ego sum, dicit Dóminus: de quacúmque tribulatióne clamáverint ad
me, exáudiam eos: et ero illórum Dóminus in perpétuum. \textnormal{Ps.\ 77, 1}
Atténdite, pópule meus, legem meam: inclináte aurem vestram in verba oris mei.
℣.~Glória Patri.
\end{latinproper}

\begin{englishwitness}{Cummiskey, The Roman Missal (1861), XIX Sunday after Pentecost, Introit}
I am the Saviour of my people, saith the Lord: in whatever distress they call
on me, I will hear them: and will be their Lord for ever.
\end{englishwitness}

\begin{englishwitness}{Douay--Rheims (Challoner), Ps 77:1, the verse}
Attend, O my people, to my law: incline your ears to the words of my mouth.
\end{englishwitness}

\noindent The antiphon is not a verse of the Vulgate, and the Missal cites no
source for it; the \latin{Ps.~77, 1} it prints stands before the verse alone.
Its phrases have the sound of the psalter's promises of salvation and hearing:
``say to my soul: I am thy salvation'' (Ps 34:3), ``He shall cry to me, and I
will hear him: I am with him in tribulation'' (Ps 90:15). None of them is its
source. The 1861 English renders \latin{Salus} as ``Saviour'' and adds ``my''
to \latin{populi}; the Latin says ``the salvation of the people''. The Vulgate
counts the psalm's title, \latin{Intellectus Asaph}, ``Understanding for
Asaph'', inside its first verse; the Missal's verse begins after it, and so
does the English printed here. After the \latin{Gloria Patri} the antiphon is
repeated, as the general rubrics direct (RGMR 427).

\subsection{Collect}
\label{proper-collect}

\properlocus{\latin{Oratio}. Missal p.~411, no.~1680.}

\begin{latinproper}
Omnípotens et miséricors Deus, univérsa nobis adversántia propitiátus exclúde:
ut mente et córpore páriter expedíti, quæ tua sunt, líberis méntibus
exsequámur. Per Dóminum.
\end{latinproper}

\begin{englishwitness}{Cummiskey, The Roman Missal (1861), XIX Sunday after Pentecost, Collect}
O Almighty and merciful God, favourably defend us from all adversity: that
being free both in soul and body, we may with security of mind perform thy
service. Thro'.
\end{englishwitness}

\noindent The Collect is said alone, with no other oration added to it. The 1861 English is
free at three points. It has God ``defend us from'' what the Latin asks him,
\latin{propitiatus}, to shut out, \latin{exclude}; its ``security of mind''
stands for \latin{liberis mentibus}, with free minds; and its ``thy service''
stands for \latin{quae tua sunt}, the things that are thine. The participle
\latin{expediti}, disencumbered, set free for action, is the word the
Postcommunion takes up at the end of the Mass. The conclusion is printed short
and is said in the full form of the general rubrics (RG 115~a).

\subsection{Epistle}
\label{proper-epistle}

\properlocus{\latin{Lectio Epistolae beati Pauli Apostoli ad Ephesios}. Eph 4:23--28. Missal p.~411, no.~1681.}

\begin{latinproper}
Fratres: Renovámini spíritu mentis vestræ, et indúite novum hóminem, qui
secúndum Deum creátus est in iustítia et sanctitáte veritátis. Propter quod
deponéntes mendácium, loquímini veritátem unusquísque cum próximo suo: quóniam
sumus ínvicem membra. Irascímini, et nolíte peccáre: sol non óccidat super
iracúndiam vestram. Nolíte locum dare diábolo: qui furabátur, iam non furétur;
magis autem labóret, operándo mánibus suis, quod bonum est, ut hábeat unde
tríbuat necessitátem patiénti.
\end{latinproper}

\begin{englishwitness}{Douay--Rheims (Challoner), Eph 4:23--28}
And be renewed in the spirit of your mind: And put on the new man, who
according to God is created in justice and holiness of truth. Wherefore,
putting away lying, speak ye the truth, every man with his neighbour. For we
are members one of another. Be angry: and sin not. Let not the sun go down upon
your anger. Give not place to the devil. He that stole, let him now steal no
more: but rather let him labour, working with his hands the thing which is
good, that he may have something to give to him that suffereth need.
\end{englishwitness}

\noindent The Missal opens the lesson with the address \latin{Fratres} and
leaves out the connective with which the Vulgate joins verse 23 to the verse
before it (\latin{Renovamini autem}); the English begins with that connective,
``And''. Otherwise the Latin agrees with the Vulgate word for word.

\subsection{Gradual}
\label{proper-gradual}

\properlocus{Ps 140:2 (Heb.\ 141), respond and verse. Missal p.~411, no.~1682.}

\begin{latinproper}
Dirigátur orátio mea sicut incénsum in conspéctu tuo, Dómine. ℣.~Elevátio
mánuum meárum sacrifícium vespertínum.
\end{latinproper}

\begin{englishwitness}{Douay--Rheims (Challoner), Ps 140:2}
Let my prayer be directed as incense in thy sight; the lifting up of my hands,
as evening sacrifice.
\end{englishwitness}

\noindent Both halves of the chant are the one verse. The Missal adds the
vocative \latin{Domine} after \latin{in conspectu tuo}, which the Vulgate's verse
does not have and the English therefore does not render.

\subsection{Alleluia}
\label{proper-alleluia}

\properlocus{Ps 104:1 (Heb.\ 105). Missal p.~411, no.~1683.}

\begin{latinproper}
Allelúia, allelúia. ℣.~\textnormal{Ps.\ 104, 1} Confitémini Dómino, et invocáte
nomen eius: annuntiáte inter gentes ópera eius. Allelúia.
\end{latinproper}

\begin{englishwitness}{Douay--Rheims (Challoner), Ps 104:1}
Give glory to the Lord, and call upon his name: declare his deeds among the
Gentiles.
\end{englishwitness}

\noindent The Vulgate counts the psalm's title, \latin{Alleluia}, inside its
first verse; the Missal's verse begins after it. There is no Tract and no
Sequence.

\subsection{Gospel}
\label{proper-gospel}

\properlocus{\latin{Sequentia sancti Evangelii secundum Matthaeum}. Mt 22:1--14. Missal pp.~411--412, no.~1684.}

\begin{latinproper}
In illo témpore: Loquebátur Iesus princípibus sacerdótum et pharisǽis in
parábolis, dicens: Símile factum est regnum cælórum hómini regi, qui fecit
núptias fílio suo. Et misit servos suos vocáre invitátos ad núptias, et
nolébant veníre. Iterum misit álios servos, dicens: Dícite invitátis: Ecce
prándium meum parávi, tauri mei et altília occísa sunt, et ómnia paráta: veníte
ad núptias. Illi autem neglexérunt: et abiérunt, álius in villam suam, álius
vero ad negotiatiónem suam: réliqui vero tenuérunt servos eius, et contuméliis
afféctos occidérunt. Rex autem cum audísset, irátus est: et missis exercítibus
suis, pérdidit homicídas illos, et civitátem illórum succéndit. Tunc ait servis
suis: Núptiæ quidem parátæ sunt, sed qui invitáti erant, non fuérunt digni. Ite
ergo ad éxitus viárum, et quoscúmque invenéritis, vocáte ad núptias. Et egréssi
servi eius in vias, congregavérunt omnes, quos invenérunt, malos et bonos: et
implétæ sunt núptiæ discumbéntium. Intrávit autem rex, ut vidéret
discumbéntes, et vidit ibi hóminem non vestítum veste nuptiáli. Et ait illi:
Amíce, quómodo huc intrásti non habens vestem nuptiálem? At ille obmútuit. Tunc
dixit rex minístris: Ligátis mánibus et pédibus eius, míttite eum in ténebras
exterióres: ibi erit fletus et stridor déntium. Multi enim sunt vocáti, pauci
vero elécti.
\end{latinproper}

\begin{englishwitness}{Douay--Rheims (Challoner), Mt 22:2--14}
The kingdom of heaven is likened to a king who made a marriage for his son. And
he sent his servants to call them that were invited to the marriage: and they
would not come. Again he sent other servants, saying: Tell them that were
invited, Behold, I have prepared my dinner: my beeves and fatlings are killed,
and all things are ready. Come ye to the marriage. But they neglected and went
their ways, one to his farm and another to his merchandise. And the rest laid
hands on his servants and, having treated them contumeliously, put them to
death. But when the king had heard of it, he was angry: and sending his armies,
he destroyed those murderers and burnt their city. Then he saith to his
servants: The marriage indeed is ready; but they that were invited were not
worthy. Go ye therefore into the highways; and as many as you shall find, call
to the marriage. And his servants going forth into the ways, gathered together
all that they found, both bad and good: and the marriage was filled with
guests. And the king went in to see the guests: and he saw there a man who had
not on a wedding garment. And he saith to him: Friend, how camest thou in
hither not having on a wedding garment? But he was silent. Then the king said
to the waiters: Bind his hands and feet, and cast him into the exterior
darkness. There shall be weeping and gnashing of teeth. For many are called,
but few are chosen.
\end{englishwitness}

\noindent The Missal's opening, \latin{In illo tempore: Loquebatur Iesus
principibus sacerdotum et pharisaeis in parabolis, dicens}, stands in place of
the Gospel's first verse, ``And Jesus answering, spoke again in parables to
them, saying''. It names those to whom the parable is spoken, ``the chief
priests and Pharisees'' who, at the close of the preceding parable, knew that
Jesus spoke of them (21:45); for that reason the English above begins at
verse 2. At verse 13
the Missal reads \latin{Tunc dixit rex}, ``Then the king said'', where the
Clementine Vulgate reads the present \latin{dicit}; the English has the past.
Every other word agrees with the Vulgate. The printed cue \latin{Credo}
follows the last line.

\subsection{Offertory}
\label{proper-offertory}

\properlocus{Ps 137:7 (Heb.\ 138). Missal p.~412, no.~1685.}

\begin{latinproper}
Si ambulávero in médio tribulatiónis, vivificábis me, Dómine: et super iram
inimicórum meórum exténdes manum tuam, et salvum me fáciet déxtera tua.
\end{latinproper}

\begin{englishwitness}{Douay--Rheims (Challoner), Ps 137:7}
If I shall walk in the midst of tribulation, thou wilt quicken me: and thou
hast stretched forth thy hand against the wrath of my enemies: and thy right
hand hath saved me.
\end{englishwitness}

\noindent The English carries the first clause and not the rest of the chant.
The Missal adds \latin{Domine} after \latin{vivificabis me}, and it sets the
last two verbs in the future, \latin{extendes} and \latin{faciet}, where the
Vulgate has the perfect, \latin{extendisti} and \latin{fecit}, and the Douay
the past, ``thou hast stretched forth'' and ``hath saved''. In the chant,
therefore, the psalmist's confidence in God's hand and right hand looks
forward, as his confidence that he will be kept alive already does in the
Vulgate.

\subsection{Secret}
\label{proper-secret}

\properlocus{\latin{Secreta}. Missal p.~412, no.~1686.}

\begin{latinproper}
Hæc múnera, quǽsumus, Dómine, quæ óculis tuæ maiestátis offérimus, salutária
nobis esse concéde. Per Dóminum nostrum Iesum Christum, Fílium tuum: Qui tecum
vivit et regnat in unitáte.
\end{latinproper}

\begin{englishwitness}{Cummiskey, The Roman Missal (1861), XIX Sunday after Pentecost, Secret}
Grant, we beseech thee, O Lord, that the offerings we bring before thy divine
Majesty may avail to our salvation. Thro'.
\end{englishwitness}

\noindent Where the English has ``before thy divine Majesty'', the Latin offers
the gifts \latin{oculis tuae maiestatis}, to the eyes of God's majesty. The
Missal prints a longer conclusion here than on the Collect and Postcommunion,
and follows the Secret with the direction \latin{Praefatio de Ss.ma
Trinitate}.

\subsection{Communion}
\label{proper-communion}

\properlocus{Ps 118:4--5 (Heb.\ 119). Missal p.~412, no.~1687.}

\begin{latinproper}
Tu mandásti mandáta tua custodíri nimis: útinam dirigántur viæ meæ, ad
custodiéndas iustificatiónes tuas.
\end{latinproper}

\begin{englishwitness}{Douay--Rheims (Challoner), Ps 118:4--5}
Thou hast commanded thy commandments to be kept most diligently. O! that my
ways may be directed to keep thy justifications.
\end{englishwitness}

\noindent The chant agrees with the Vulgate word for word.

\subsection{Postcommunion}
\label{proper-postcommunion}

\properlocus{\latin{Postcommunio}. Missal p.~412, no.~1688.}

\begin{latinproper}
Tua nos, Dómine, medicinális operátio, et a nostris perversitátibus cleménter
expédiat, et tuis semper fáciat inhærére mandátis. Per Dóminum.
\end{latinproper}

\begin{englishwitness}{Cummiskey, The Roman Missal (1861), XIX Sunday after Pentecost, Postcommunion}
May the healing efficacy of these thy mysteries, O Lord, mercifully free us from
our perverseness, and make us always obedient to thy commandments. Thro'.
\end{englishwitness}

\noindent ``These thy mysteries'' is the translator's addition: the Latin names
only \latin{tua medicinalis operatio}, thy healing working. ``Obedient'' stands
for \latin{inhaerere}, to cleave or cling. The verb \latin{expediat} is the
Collect's \latin{expediti} again, so that the Mass asks at its end for the
freedom it asked for at its beginning.
